% Chapter 15 -- Care and Maintenance (the 1991 Macintosh User's Guide Ch. 15
% "Care and Maintenance", apple-macintosh.md Sec 8.2).  Sources (Law 1):
% os/milton/kmain.c (Shut Down / Restart / Exit to InitechDOS commit
% DESKTOP.DB; TRASH-CREATE and DESKTOP-DB-CREATE when they are missing);
% os/flair/finder_ops.c (the Trash is a folder on the disk until emptied);
% os/apps/itext/itext.c (unsaved work is in memory only); the InitechDOS
% books for CHKDSK, COPY and XCOPY.  The hardware advice is general care of
% the period's equipment and makes no claim about the software.
% GATES: test-flair-shutdown test-desktop-db test-flair-empty-trash test-flair-itext
\Gates{test-flair-shutdown, test-desktop-db, test-flair-empty-trash,
  test-flair-itext}
\chapter{Care and Maintenance}
\label{ch:care}
\index{care and maintenance}

\section{Protecting your work}
\SideHead{Save}A document is in memory until you save it.  Save it every
few minutes, and before you leave the computer.
\SideHead{Shut down}Always use \menuitem{Shut Down} before you turn the
computer off.  It saves the arrangement of the desktop and gives the open
application the chance to save your work.
\SideHead{Empty the Trash with care}Items in the Trash still take up space
on the disk; emptying the Trash frees it, and cannot be undone.
\SideHead{Keep copies}Copy your documents to another disk at the end of the
day, with \file{COPY} or \file{XCOPY} at the InitechDOS prompt
(Chapter~\ref{ch:dos}).

\section{Checking the data disk}
\index{CHKDSK}
From time to time, check the data disk with \file{CHKDSK}: leave FLAIR,
change to the \file{APPS} directory and type \type{CHKDSK} (Chapter~\ref{ch:dos}).
The \doctitle{InitechDOS Reference} explains its report and what to do if
it finds a problem.

\section{When the disk is nearly full}
A new folder cannot be made, and documents cannot be saved, when the data
disk is full.  Empty the Trash, then move documents you no longer need
onto another disk at the InitechDOS prompt.

\section{Looking after the equipment}
\begin{itemize}
\item Keep the computer, the display and the disks away from heat, damp
  and dust, and leave room round the computer for air to flow.
\item Keep disks away from magnets, telephones and the display, and do not
  touch their surfaces.
\item Clean the display with a soft cloth; never spray liquid on it.
\item Turn the computer off before you connect or disconnect the keyboard,
  the mouse or the display.
\item Clean the mouse ball and rollers when the pointer stops moving
  smoothly.
\end{itemize}
